% Edit the title, individual panel descriptions, and shared notes here.
% All figure placements read these commands; edit once to update every continuation.
\newcommand{\FigThreeTitle}{TrajRisk architecture, training diagnostics, and scale-specific representations.}
\newcommand{\FigThreeA}{Four parallel scale branches each apply two directed, edge-weighted TrajRiskConv layers to 31-dimensional node features. Each branch produces a 64-dimensional embedding; the concatenated 256-dimensional representation is passed to a fusion multilayer perceptron (MLP) to predict a dimensionless TrajRisk trajectory-state score between 0 and 1. The convolution combines learned neighbor and self transforms without additional symmetric graph convolutional network (GCN) degree normalization.}
\newcommand{\FigThreeB}{Training and validation mean squared error (MSE) across epochs. The lowest validation MSE occurs at epoch 119.}
\newcommand{\FigThreeC}{Validation AUROC across epochs. The AUROC-best checkpoint at epoch 144 (AUROC = 0.9071) is the restored model used for downstream analyses; it is distinct from the MSE-best checkpoint. Dashed vertical lines mark learning-rate reductions.}
\newcommand{\FigThreeD}{Validation Spearman rank correlation between predicted TrajRisk score and normalized pseudotime, reaching \ensuremath{\rho} = 0.5814 at epoch 144.}
\newcommand{\FigThreeE}{TrajRisk-score distributions for all 6,582 Walesky hepatocytes, with sampling-time counts shown below the axis. White circles indicate medians, thick black segments span the interquartile range, and numerical annotations indicate group means.}
\newcommand{\FigThreeF}{Spearman correlation between the \ensuremath{L_2} norm of each cell's 64-dimensional scale embedding and its final predicted TrajRisk score, across the Walesky cohort. This is an association summary rather than a causal measure of scale importance.}
\newcommand{\FigThreeNotes}{The fixed timepoint-stratified split contains 5,264 training and 1,318 validation cells. Plotted ``risk'' denotes the model estimate of normalized pseudotime, not an outcome probability. Normalized pseudotime is the regression target; binary labels for AUROC evaluation use pseudotime\_norm > 0.10. AUROC, area under the receiver operating characteristic curve. Scales 1--4 connect 0--6, 6--24, 24--48, and 48--96~h, respectively.}
